\documentclass[11pt]{article}
\usepackage{fontspec}
\usepackage[french]{babel}
\usepackage{amsmath}
\usepackage{amssymb}
\usepackage{amsthm}
\usepackage{mathrsfs}
\usepackage{tikz-cd}
\usepackage{enumitem}

\frenchsetup{AutoSpacePunctuation=false}

\newcommand{\printedpage}[1]{\marginpar{\footnotesize p.~#1}}
\newcommand{\msp}[1]{\marginpar{\footnotesize [#1]}}

\begin{document}

\printedpage{120}
% head: A. Grothendieck : Dérivateurs, Ch. XII, Édition des Universités de Montpellier II et Paris VII

\msp{page 180}
$\mathrm{b}_n \Longrightarrow \mathrm{c}_n$. On peut supposer que $f$ est une Hot-fibration (et même Hot-parfait si ça nous chante \dots). On sait déjà que pour
\[
H^i(Y,F) \longrightarrow H^i(X,f^*F) \tag*{$(*)$}
\]
tout est O.K. si $i \in \{0,1\}$. Donc on peut supposer $F$ abélien. Si ça nous avance, on peut supposer, par récurrence, (que l’implication $\mathrm{b}_{n-1} \Longrightarrow \mathrm{c}_{n-1}$ est connue, donc) que $(*)$ est \texttt{[un]} isomorphisme pour $i<n$. Reste à le voir pour $i=n$. On est ramené bien sûr, comme $\pi_0(f)$ est bijectif, au cas où $X$, $Y$ sont 0-connexes, donc il y a essentiellement une seule fibre $Z$. La suite exacte d’homotopie nous apprend que
\[
\pi_i(Z,z)=0 \qquad \text{pour } i \le n-1 \ (Z \text{ est } (n-1)\text{-connexe})
\]
et de plus
\[
\pi_n(Z,z) \longrightarrow \pi_n(X,z) \qquad \text{est nul}
\]
(car $\pi_n(X,z) \longrightarrow \pi_n(Y,y)$ \texttt{[est un]} isomorphisme donc \texttt{[un]} monomorphisme), donc
\[
\pi_{n+1}(Y,y) \longrightarrow \pi_n(Z,z) \qquad \text{\texttt{[est un]} épimorphisme.}
\]
Comme $Z$ est $(n-1)$-connexe (donc aussi 1-connexe), on a
\[
H^i(Z,G)=0\ , \qquad \text{si } 0<i<n\ , \tag*{$(*)$}
\]
pour $G$ localement constant (donc constant) sur $Z$.

\msp{page 181}
Écrivons la suite spectrale de Leray pour $f : X \longrightarrow Y$ et $f^*F = F_X$ :
\[
H^*(X,f^*F) \Longleftarrow E_2^{pq} = H^p(Y,H^q(Z,f^*F)),
\]
où ‘$H^q(Z)$’ doit être vu comme un système local sur $Y$. Donc par $(*)$ on a
\[
E_2^{pq}=0 \qquad \text{si } 0<q<n, \qquad
\begin{tikzpicture}[baseline=0pt, x=0.6cm, y=0.6cm]
\draw[->] (0,0) -- (0,4);
\draw[->] (0,0) -- (6,0);
\node[left] at (0,3) {$n$};
\foreach \i in {0,...,4} { \fill (\i,3) circle (2.5pt); \fill (\i,0) circle (2.5pt); }
\node at (1,1.5) {$0$};
\node at (3,1.5) {$0$};
\end{tikzpicture}
\qquad ,
\]
donc on trouve la suite exacte
\[
0 \longrightarrow H^n(Y,F) \longrightarrow H^n(X,F_X) \xrightarrow{\ \alpha\ } H^0(Y,\ \underbrace{H^n(Z,\overbrace{F_Z}^{M_Z})}_{\begin{array}{l}\scriptstyle = \ \mathrm{Hom}(\pi_n(Z,z),M)\\ \scriptstyle \text{par \texttt{[le]} th. de Hurewicz}\end{array}}\ ) \longrightarrow H^{n+1}(Y,F).
\]

\printedpage{121}
% head: A. Grothendieck : Dérivateurs, Ch. XII, Édition des Universités de Montpellier II et Paris VII

Ici je fixe un point $x$ de $X$, $Z$ désigne la fibre de $f$ en $y=f(x)$, et $M$ celle de $F$ en $y$, $M=F_y$. Le groupe $\pi_1(Y)$ y opère, d’ailleurs $\pi \overset{\text{déf}}{=} \pi_1(X,x) \simeq \pi_1(Y,y)$ aussi. Le $H^0(Y,H^n(Z,F_Z))$ peut s’interpréter, sauf erreur, comme
\[
\mathrm{Hom}_\pi(\pi_n(Z,x),M),
\]
en tenant compte de l’opération de $\pi=\pi_1(X,x)$ sur $\pi_n(Z,x)$ également. (À vérifier.) Il faut interpréter la flèche
\[
H^n(X,F_X) \longrightarrow \mathrm{Hom}_\pi(\pi_n(Z,x),M),
\]
via l’application composée \texttt{[où $z=x$, vu comme point de $Z$]}
\[
\pi_n(Z,z) \longrightarrow \pi_n(X,x) \longrightarrow H_n(X,\mathbf{Z}). \tag*{$(*)$}
\]
Si $F$ est \textit{constant} sur $Y$ (p. ex. si $Y$ \texttt{[est]} 1-connexe), donc de valeur $M$, on a les flèches canoniques
\[
H^n(X,F_X) = H^n(X,M) \longrightarrow \mathrm{Hom}(H_n(X,\mathbf{Z}),M) \longrightarrow \mathrm{Hom}(\pi_n(Z,z),M)\text{(\footnotemark[118])} \tag*{$(*')$}
\]
\footnotetext[118]{grâce à la composition avec $(*)$.}
\texttt{[lire : }$\mathrm{Hom}_\pi$\texttt{]},

\msp{page 182}
et sûrement (en dehors de toute hypothèse de type $\mathrm{b}_n$ ou autre) l’application canonique directe ‘de localisation sur $Y$’, suivi de \dots,
\[
H^n(X,F_X) \longrightarrow H^0(Y,H^n(Z,F_Z)) \longrightarrow \mathrm{Hom}_\pi(\pi_n(Z),M) \tag*{$(**)$}
\]
n’est autre que la composée des applications canoniques ci-dessus. Dans le cas actuel où on suppose $\mathrm{b}_n$, la deuxième flèche dans $(**)$ est \texttt{[un]} isomorphisme, donc la composée ‘n’est autre’ que la flèche de localisation elle-même. D’autre part $\pi_n(Z,z) \longrightarrow \pi_n(X,x)$ dans $(*)$ est nul, donc aussi $\pi_n(Z,z) \longrightarrow H_n(X,\mathbf{Z})$, donc aussi la flèche composée $(*')$ qui s’en déduit. Revenons alors à la suite exacte de Leray (page précédente), on trouve que $H^n(Y,F) \xrightarrow{\ \sim\ } H^n(X,F_X)$.

Comment récupérer le cas où $F$ n’est \textit{pas} constant? Dans le cas $n=2$, cela marchait, en passant aux revêtements universels $\tilde{X}$, $\tilde{Y}$ de $X$, $Y$ (par rapport aux points $x$, $y$ choisis), et en écrivant les suites spectrales
\[
\begin{tikzcd}[row sep=small, column sep=small]
H^*(Y,F) \arrow[d] & \Longleftarrow & E_2^{pq} \arrow[d] & = & H^p(\pi,H^q(\tilde{Y},F_{\tilde{Y}})) \\
H^*(X,F_X) & \Longleftarrow & E^{\prime pq}_2 & = & H^p(\pi,H^q(\tilde{X},F_{\tilde{X}})).
\end{tikzcd}
\]
L’hypothèse $\mathrm{b}_n$ est conservée en passant de $f : X \longrightarrow Y$ à $\tilde{f} = \tilde{f}_x : \tilde{X} \longrightarrow \tilde{Y}$, donc on sait que
\printedpage{122}
% head: A. Grothendieck : Dérivateurs, Ch. XII, Édition des Universités de Montpellier II et Paris VII
\msp{page 183}
\[
E_2^{pq} \xrightarrow{\ \sim\ } E^{\prime pq}_2 \qquad \text{\texttt{[est un]} isomorphisme si } q \le n. \tag*{$(*)$}
\]
Dans le cas $n=2$, utilisant que
\[
E_2^{p1} = E^{\prime p1}_2 = 0 \qquad \text{car } H^1(\tilde{X}) = H^1(\tilde{Y}) = 0, \qquad
\begin{tikzpicture}[baseline=0pt, x=0.5cm, y=0.5cm]
\draw[->] (0,0) -- (0,2.8);
\draw[->] (0,0) -- (5,0);
\node[above left] at (0,2.6) {$q$};
\node[left] at (0,2) {$2$};
\foreach \i in {0,...,3} { \fill (\i,2) circle (2.5pt); \fill (\i,0) circle (2.5pt); }
\node at (0,1) {$0$};
\node at (1,1) {$0$};
\node at (2,1) {$0$};
\end{tikzpicture}
\qquad ,
\]
on trouve deux suites exactes
\[
\begin{tikzcd}
0 \arrow[r] & H^2(\pi,H^0(\tilde{Y})) \arrow[r] \arrow[d, "\text{iso}"] & H^2(Y) \arrow[r] \arrow[d] & H^0(\pi,H^2(\tilde{Y})) \arrow[r] \arrow[d, "\text{iso}"] & H^3(\pi,H^0(\tilde{Y})) \arrow[d, "\text{iso}"] \\
0 \arrow[r] & H^2(\pi,H^0(\tilde{X})) \arrow[r] & H^2(X) \arrow[r] & H^0(\pi,H^2(\tilde{X})) \arrow[r] & H^3(\pi,H^0(\tilde{X})),
\end{tikzcd}
\]
et un homomorphisme de l’une dans l’autre, qui permet de conclure par le lemme des 5 que $H^2(Y,F) \longrightarrow H^2(X,F_X)$ est \texttt{[un]} isomorphisme. Mais pour $n$ quelconque, la relation $(*)$ sur les termes $E_2$ (pour $q \le n$) ne me semble pas impliquer la relation similaire pour les $E_n$ supérieures, et de là pour les termes $E_\infty$, ce qui permettrait de conclure.

Je reviens donc à la suite exacte p. 181, avec $F$ faisceau \textit{tordu.} J’ai envie de trouver une application canonique
\[
H^n(X,G) \longrightarrow \overbrace{\mathrm{Hom}_\pi(\pi_n(X,x),G_x)}^{\text{\scriptsize Homs \textit{additifs} bien sûr, i.e. de $\mathbf{Z}[\pi]$-modules}} \qquad ,
\]
où $G$ est un faisceau abélien localement constant quelconque sur $X$, et $\pi = \pi_1(X,x)$. Si j’interprète
\msp{page 184}
$\pi_n(X,x)$ par des classes d’homotopie d’applications ponctuées
\[
(S^n,x_0) \xrightarrow{\ \varphi\ } (X,x)
\]
j’ai bien, pour une telle application $\varphi$,
\[
\varphi^* : H^n(X,G) \longrightarrow H^n(S^n,\ \overbrace{\varphi^*G}^{\simeq\,(G_x)_{S^n}}\ ) \simeq G_x
\]
(tenant compte que $S^n$ est 1-connexe, et que sa $n$-cohomologie à valeurs dans un faisceau constant $M$ est $M$.) Sûrement l’application
\[
H^n(X,G) \longrightarrow G_x
\]
\printedpage{123}
% head: A. Grothendieck : Dérivateurs, Ch. XII, Édition des Universités de Montpellier II et Paris VII
ainsi définie ne dépend que de la classe d’homotopie ponctuée de $\varphi$, d’où
\[
\pi_n(X,x) \times H^n(X,G) \longrightarrow G_x,
\]
et il faut de plus prouver a) que c’est compatible avec les opérations de $\pi = \pi_1(X,x)$, et b) que c’est additif en $\pi_n(X,x)$, pour avoir l’application cherchée
\[
H^n(X,G) \longrightarrow \mathrm{Hom}_\pi(\pi_n(X,x),G_x).
\]
Si j’admets ceci, il n’y a qu’un pas pour admettre aussi que la flèche qui nous intéresse vraiment (dans la suite exacte p. 181)
\msp{page 185}
\[
H^n(X,F_X) \longrightarrow \mathrm{Hom}_\pi(\pi_n(Z,z),\underbrace{M}_{=\,F_y})
\]
n’est autre que la composée de la flèche précédente (page précédente) pour $G = F_X$, avec
\[
\mathrm{Hom}_\pi(\pi_n(X,x),M) \longrightarrow \mathrm{Hom}_\pi(\pi_n(Z,x),M)
\]
déduite de
\[
\pi_n(Z,x) \longrightarrow \pi_n(X,x).
\]
Or cette flèche est nulle, par l’hypothèse $\mathrm{b}_n$. Donc la flèche $\alpha$ dans la suite exacte page 181 est nulle, donc $H^n(Y,F) \longrightarrow H^n(X,F_X)$ est \texttt{[un]} isomorphisme, q.e.d. (mais il y a pas mal de travail de vérification à faire \dots).

$\mathrm{c}_n \Longrightarrow \mathrm{a}_{n-1}$. On peut procéder par récurrence sur $n$. (NB L’enoncé a un sens même pour $n=1$, où il est prouvé, et même $\mathrm{c}_1 \Longrightarrow \mathrm{b}_1$, et $\mathrm{b}_1 \Longrightarrow \mathrm{a}_0$, plus généralement $\mathrm{b}_n \Longrightarrow \mathrm{a}_{n-1}$, immédiat par \texttt{[la]} suite exacte d’homotopie.) Donc on sait déjà, en plus de l’hypothèse $\mathrm{c}_n$, qu’on a $\mathrm{a}_{n-2}$, i.e. $Z$ est $(n-2)$-connexe. Il faut prouver qu’il est même $(n-1)$-connexe, i.e. $\pi_{n-1}(Z)=0$. Le cas $n=2$ demande % [sic] L’enoncé -> L’énoncé, missing accent on é
\msp{page 186}
peut-être une attention particulière, car on ne sait pas même encore que $Z$ est 1-connexe. Mais on sait que $\mathrm{c}_2$ implique $\mathrm{c}_1$, qui équivaut \texttt{[?]} à $\mathrm{b}_1$, donc on a
\[
\pi_1(X,x) \longrightarrow \pi_1(Y,y) \qquad \text{isomorphisme,}
\]
donc
\[
\pi_1(Z,x) \longrightarrow \pi_1(X,x) \qquad \text{\texttt{[est]} \textit{nul,}}
\]
donc
\[
\pi_2(Y,y) \longrightarrow \pi_1(Z,x) \qquad \text{\texttt{[est un]} épimorphisme,}
\]
ce qui implique tout au moins que $\pi_1(Z,x)$ est \textit{abélien.} Il faut le voir comme un $\pi$-module, et prouver sa nullité comme tel (moyennant $\mathrm{c}_2$).

\printedpage{124}
% head: A. Grothendieck : Dérivateurs, Ch. XII, Édition des Universités de Montpellier II et Paris VII

Passons au cas général. On a pour $F$ localement constant abélien sur $Y$
\[
H^*(X,F_X) \Longleftarrow \underbrace{E_2^{pq} = H^p(Y,H^q(Z,F_Z))}_{\begin{array}{l}\scriptstyle = 0 \text{ si } 0<q<n-1\\ \scriptstyle \text{car } Z \text{ est } (n-2)\text{-connexe}\end{array}}, \qquad
\begin{tikzpicture}[baseline=0pt, x=0.6cm, y=0.6cm]
\draw[->] (0,0) -- (0,5);
\draw[->] (0,0) -- (6,0);
\node[left] at (0,4) {$n-1$};
\foreach \i in {0,...,4} { \fill (\i,4) circle (2.5pt); \fill (\i,0) circle (2.5pt); }
\node at (1,2) {$0$};
\node at (3,2) {$0$};
\end{tikzpicture}
\qquad ,
\]
D’où
\[
0 \longrightarrow H^{n-1}(Y,F) \longrightarrow H^{n-1}(X,F_X) \longrightarrow \underbrace{H^0(Y,H^{n-1}(Z,F_Z))}_{\begin{array}{l}\scriptstyle \mathrm{Hom}_\pi(\pi_{n-1}(Z,x),F_x)\\ \scriptstyle \text{par Hurewicz}\end{array}} \longrightarrow H^n(Y,F) \longrightarrow H^n(X,F_X).
\]
L’hypothèse $\mathrm{c}_n$ implique qu’on a
\msp{page 187}
\[
\mathrm{Hom}_\pi(\pi_{n-1}(Z,x),\underbrace{M}_{=\,F_x}) = 0
\]
et ceci \textit{quel que soit le $\pi$-module $M$,} puisque $\pi = \pi_1(X,x) \xrightarrow{\ \sim\ } \pi_1(Y,y)$. Cela implique (en prenant p. ex. $M = \pi_{n-1}(Z,x)$ !) que $\pi_{n-1}(Z,x)=0$. (Mais en faisant attention, dans le cas $n=2$, qu’on sait déjà que $\pi_1(Z,x)$ est abélien).

Donc le théorème rectifie page 178 a bien l’air vrai. Il est un peu triste, puisque dans chacune des implications % [sic] rectifie -> rectifié, past participle intended
\[
\mathrm{a}_n \ \Longrightarrow\ \mathrm{b}_n \ \Longrightarrow\ \mathrm{c}_n \ \Longrightarrow\ \mathrm{a}_{n-1}
\]
on perd un chouia, et en fin de compte on sort avec la magnifique implication
\[
\mathrm{a}_n \ \Longrightarrow\ \mathrm{a}_{n-1}.
\]
De toutes façons, il est $\pm$ clair que le “bon” localiseur $W_n$, celui qui doit définir les “$n$-types d’homotopie” comme objets de $\mathrm{Hot}_{W_n}$, est défini par la condition $\mathrm{a}_n$.

\section*{2\quad Critère de $W_n$-asphéricité.}

Je vais essayer de faire mieux quand $Y$ est réduit à un point. Cette fois, je

\end{document}
